%%%Début de la configuration de xkeyval%%%
\makeatletter
%%%%%%%%%%%%%%%%%%%Pour les sujets%%%%%%%%%%%%%%%%

%% 1. Définition des clés pour les sujets 
\define@cmdkey	[SUJ]	{sujet}	{semestre}		{}
\define@cmdkey	[SUJ]	{sujet}	{lieu}			{}
\define@cmdkey	[SUJ]	{sujet}	{annee}			{}
\define@cmdkey	[SUJ]	{sujet}	{type}			{}
\define@cmdkey	[SUJ]	{sujet}	{jour}			{}
\define@cmdkey	[SUJ]	{sujet}	{prenomauteur}	{}
\define@cmdkey	[SUJ]	{sujet}	{initialeprenom}{}
\define@cmdkey	[SUJ]	{sujet}	{nomauteur}		{}
\define@cmdkey	[SUJ]	{sujet}	{titre}			{}
\define@cmdkey	[SUJ]	{sujet}	{datepub}		{}
\define@cmdkey	[SUJ]	{sujet}	{ajoutref}		{}
\define@cmdkey	[SUJ]	{sujet}	{genre}			{}
\define@cmdkey	[SUJ]	{sujet}	{corrige}		{}
\define@cmdkey	[SUJ]	{sujet}	{liencorrige}	{}
%\define@cmdkey	[SUJ]	{sujet}	{texte}			{}
\define@cmdkey	[SUJ]	{sujet}	{typequn}		{}
\define@cmdkey	[SUJ]	{sujet}	{typeqdeux}		{}
\define@cmdkey	[SUJ]	{sujet}	{qun}			{}
\define@cmdkey	[SUJ]	{sujet}	{qdeux}			{}
\define@cmdkey	[SUJ]	{sujet}	{themeprogrammequn}	{}
\define@cmdkey	[SUJ]	{sujet}	{themeprogrammeqdeux}	{}

%% 2. Définition des commandes pour accéder aux valeurs des clés
\newcommand{\semestre}{\cmdSUJ@sujet@semestre}
\newcommand{\lieu}{\cmdSUJ@sujet@lieu}
\newcommand{\annee}{\cmdSUJ@sujet@annee}
\newcommand{\type}{\cmdSUJ@sujet@type}
\newcommand{\jour}{\cmdSUJ@sujet@jour}
\newcommand{\prenomauteur}{\cmdSUJ@sujet@prenomauteur}
\newcommand{\initialeprenom}{\cmdSUJ@sujet@initialeprenom}
\newcommand{\nomauteur}{\cmdSUJ@sujet@nomauteur}
\newcommand{\titre}{\cmdSUJ@sujet@titre}
\newcommand{\datepub}{\cmdSUJ@sujet@datepub}
\newcommand{\ajoutref}{\cmdSUJ@sujet@ajoutref}
\newcommand{\genre}{\cmdSUJ@sujet@genre}
\newcommand{\corrige}{\cmdSUJ@sujet@corrige}
\newcommand{\liencorrige}{\cmdSUJ@sujet@liencorrige}
%\newcommand{\texte}{\cmdSUJ@sujet@texte}
\newcommand{\typequn}{\cmdSUJ@sujet@typequn}
\newcommand{\typeqdeux}{\cmdSUJ@sujet@typeqdeux}
\newcommand{\qun}{\cmdSUJ@sujet@qun}
\newcommand{\qdeux}{\cmdSUJ@sujet@qdeux}
\newcommand{\themeprogrammequn}{\cmdSUJ@sujet@themeprogrammequn}
\newcommand{\themeprogrammeqdeux}{\cmdSUJ@sujet@themeprogrammeqdeux}

%% 3. Valeurs par défaut des clés 
\presetkeys[SUJ]{sujet}{%
				nomauteur = Anonyme,
				prenomauteur =,
				initialeprenom =,
				titre = {LeTitre},
				semestre = HQRS,
				lieu = Métropole,
				annee = 2020,
				type =,
				jour = 0,
				datepub = 4545,
				ajoutref =,
				genre = \genrenondefini{},
				corrige =,
				liencorrige=,
				typequn = \litt,
				typeqdeux = \phil,
				qun =,
				qdeux =,
				themeprogrammequn=,
				themeprogrammeqdeux=,}{}
%				texte = {Lorem ipsum dolor sit amet},}
%%%Fin de la configuration de xkeyval

%% 4. Macros utilitaires pour le formatage et l'affichage
% 4.1 Formatage du nom de l'auteur selon les cas
\newcommand{\formatauteur}{%
    \ifthenelse{\equal{\prenomauteur}{}}%
    {{\textsc{\nomauteur}}}% Sans prénom
    {%
        \IfEndWith{\initialeprenom}{de}% Avec particule
        {{\initialeprenom{}~\textsc{\nomauteur}}}%
        {{\initialeprenom{}.~\textsc{\nomauteur}}}%
    }%
}

\newcommand{\sujetn}[2]{
\stepcounter{nosujet}
%%%%%Pour la mise en page%%%%
\renewcommand{\headrulewidth}{0.4pt}
\fancyhead{}
\fancyfoot{} %clear all footer fields
\fancyfoot[RO]{\textbf{\thepage ~sur~\pageref*{TotPages}}\hspace*{1.5cm}}
\fancyfoot[LE]{\hspace*{1.5cm}\textbf{\thepage ~sur~\pageref*{TotPages}}}
  \fancyfoot[LO,RE]{\scriptsize{Humanités, Littérature, Philosophie -- Banque des sujets indexés -- Terminale}}
    \setkeys[SUJ]{sujet}{#1}
    \ifthenelse{\equal{\type}{}}
    {%%%Cas où le type d'épreuve n'est pas précisé = elle est obligatoire
    \ifthenelse{\equal{\prenomauteur}{}}
    {%%%L'auteur n'a pas de prénom
	\fancyhead[C]{\lieu{}~\annee{} (Jour~\jour{})~:~\textsc{\nomauteur}, \textit{\titre}}
	\phantomsection\addcontentsline{toc}{section}{\textsc{\nomauteur}, \textit{\titre} (\lieu~\annee, Jour~\jour{})}
		\index[n]{{\annee~\lieu~(Jour~\jour) -- \nomauteur}@\textbf{\annee,~\lieu} (Jour~\jour) -- \\*\hspace*{1em}\nomauteur, \textit{\titre}}
	}
	{%%%L'auteur a un prénom
	\IfEndWith{\initialeprenom}{de}
		{%%%L'auteur a une particule dans son nom de famille
		\fancyhead[C]{\lieu{}~\annee{} (Jour~\jour{})~:~\textsc{\initialeprenom{}}~\textsc{\nomauteur{}}, \textit{\titre{}}}
\phantomsection\addcontentsline{toc}{section}{\textsc{\initialeprenom{}}~\textsc{\nomauteur}, \textit{\titre} (\lieu~\annee, Jour~\jour{})}
		\index[n]{{\annee~\lieu~(Jour~\jour) -- \nomauteur}@\textbf{\annee,~\lieu} (Jour~\jour) -- \\*\hspace*{1em} \initialeprenom{}~\nomauteur, \textit{\titre}}
		}%%%Fin du cas avec particule
{\fancyhead[C]{\lieu{}~\annee{} (Jour~\jour{})~:~\initialeprenom{}.~\textsc{\nomauteur{}}, \textit{\titre{}}}
\phantomsection\addcontentsline{toc}{section}{\initialeprenom{}.~\textsc{\nomauteur}, \textit{\titre} (\lieu~\annee, Jour~\jour{})}
		\index[n]{{\annee~\lieu~(Jour~\jour) -- \nomauteur}@\textbf{\annee,~\lieu} (Jour~\jour) -- \\*\hspace*{1em}\initialeprenom{}.~\nomauteur, \textit{\titre}}
		}%%%Fin du cas normal sans particule
	}%%%Fin du cas de l'auteur avec un prénom
	}%%%Fin du cas de l'épreuve obligatoire
	{%%%Le type d'épreuve est précisé
	 \ifthenelse{\equal{\prenomauteur}{}}
    {%%%L'auteur n'a pas de prénom
	\fancyhead[C]{\lieu{}~\annee{} (\type, Jour~\jour{})~:~\textsc{\nomauteur}, \textit{\titre}}
	\phantomsection\addcontentsline{toc}{section}{\textsc{\nomauteur}, \textit{\titre} (\lieu~\annee~(\type), Jour~\jour{})}
	\index[n]{{\annee~\lieu~(\type, Jour~\jour) -- \nomauteur}@\textbf{\annee,~\lieu}, (\type, Jour~\jour) -- \\*\hspace*{1em}\nomauteur, \textit{\titre}}
	}
	{%%%L'auteur a un prénom
	\IfEndWith{\initialeprenom}{de}
		{%%%L'auteur a une particule dans son nom de famille
		\fancyhead[C]{\lieu{}~\annee{} (\type, Jour~\jour{})~:~\textsc{\initialeprenom{}}~\textsc{\nomauteur{}}, \textit{\titre{}}}
\phantomsection\addcontentsline{toc}{section}{\textsc{\initialeprenom{}}~\textsc{\nomauteur}, \textit{\titre} (\lieu~\annee~(\type), Jour~\jour{})}
		\index[n]{{\annee~\lieu~(\type, Jour~\jour) -- \nomauteur}@\textbf{\annee,~\lieu} (\type, Jour~\jour) -- \\*\hspace*{1em}{\initialeprenom{}}.~\nomauteur, \textit{\titre}}
		}%%%Fin du cas avec particule
{\fancyhead[C]{\lieu{}~\annee{} (\type, Jour~\jour{})~:~\initialeprenom{}.~\textsc{\nomauteur{}}, \textit{\titre{}}}
\phantomsection\addcontentsline{toc}{section}{\initialeprenom{}.~\textsc{\nomauteur}, \textit{\titre} (\lieu~\annee~(\type), Jour~\jour{})}
		\index[n]{{\annee~\lieu~(\type, Jour~\jour) -- \nomauteur}@\textbf{\annee,~\lieu} (\type, Jour~\jour) -- \\*\hspace*{1em}\initialeprenom{}.~\nomauteur, \textit{\titre}}
		}%%%Fin du cas normal sans particule
	}%%%Fin du cas d'un auteur sans prénom
	}%%%Fin du cas dans lequel le type d'épreuve est précisé
	%%%Fin des test sur le type d'épreuve et le prénom
\renewcommand{\arraystretch}{1.5}
\vspace{-0.5cm}
\begin{center}
\begin{huge}
\textit{Sujet d'Humanités, Littérature, Philosophie}
\end{huge}\\
{\Large{Épreuve de Terminale}}\par
\smallskip
\hspace*{0.7cm}\noindent\fbox{%%Mise en page normal dans un petit encart au début de la page
\makebox[\linewidth-3.25cm][c]{
\begin{minipage}{\linewidth-3.25cm}
		\begin{small}
		\noindent\small \textbf{Axe du programme~:}~\semestre \hfill{
		\noindent\textbf{Genre du texte~:}~\genre}\par 
		\noindent\textbf{Auteur~:}
		\ifthenelse{\equal{\prenomauteur}{}}%%%Si l'auteur n'a pas de prénom
		{\textsc{\nomauteur}}%
		%%%Si l'auteur a un prénom
		{\IfEndWith{\initialeprenom}{de}%%%Test : il a une particule
		{\initialeprenom{}~\textsc{\nomauteur}}
		{\initialeprenom{}.~\textsc{\nomauteur}}
		}%%%Fin du cas où l'auteur a un prénom
		\hfill \textbf{Ouvrage~:} \textit{\titre}\par
		\textbf{Centre d'examen :} \lieu \hfill \textbf{Année :} \annee\par
		\textbf{Session~:~}%
		\ifthenelse{\equal{\type}{}}{Obligatoire, Jour~\jour}{\type, Jour~\jour}%, Sujet \no{}\arabic{nosujet}%Pour vérifier que le nombre de sujets est bien conforme
		\ifthenelse{\equal{\corrige}{}}{\hfill{}\textbf{Éléments de correction disponibles~:} Non}{\hfill{}\textbf{Éléments de correction disponibles~:} {\hyperlink{corr\liencorrige}{Oui}\hypertarget{\liencorrige}{}}}\end{small}
\end{minipage}
	}
	\ifthenelse{\equal{\prenomauteur}{}}
	{\index[a]{\textsc{\nomauteur}}}
	{\index[a]{\textsc{\nomauteur}, \prenomauteur}}
	}%%%Fin de la fbox
	\newline \smallskip \newline \noindent\textbf{Le candidat traite les deux parties sur des copies séparées.}
	\end{center}\vspace{-0.3cm}
	\begin{envtexte}
	#2\par
	\vspace{-0.4cm}\begin{flushright}\nopagebreak
	\ifthenelse{\equal{\prenomauteur}{}}{\textsc{\nomauteur}}
	{\prenomauteur~\textsc{\nomauteur}}, \textit{\titre} (\datepub)%
	\ifthenelse{\equal{\ajoutref}{}}{.}{, \ajoutref{}.}
	\end{flushright}
	\end{envtexte}
	\begin{envquestion}
	\noindent\question{Question d'interprétation \typequn~:}\par\nopagebreak
	\noindent\qun\index[q]{Question d'interprétation \typequn!\qun}
	\spcsujet
	\noindent\question{Essai \typeqdeux~:}\par\nopagebreak
	\noindent\qdeux\index[q]{Essai \typeqdeux!\qdeux}
	\stepcounter{sujetint\themeprogrammequn}
	\stepcounter{sujetref\themeprogrammeqdeux}
	\ifthenelse{\equal{\typequn}{\litt}}%
	%%%Si l'interprétation est littéraire, on augmente l'interprétation litt. liés et la réf. philosophique ; ainsi que les compteurs liés.
	{\stepcounter{sujetintlitteraire\themeprogrammequn}
	\stepcounter{nosujetintlitteraire}
	\stepcounter{sujetrefphilosophique\themeprogrammeqdeux}
	\stepcounter{nosujetrefphilosophique}%
	}
	{%%%Sinon, on augmente l'interprétation phil. et la réf. litt. ; ainsi que les compteurs liés.
	\stepcounter{sujetint\typequn\themeprogrammequn}
	\stepcounter{nosujetintphilosophique}
	\stepcounter{sujetreflitteraire\themeprogrammequn}
	\stepcounter{nosujetreflitteraire}
	}
%	\index[p]{{\csname\themeprogrammequn\endcsname}!\qun}
%	\index[p]{{\csname\themeprogrammeqdeux\endcsname}!\qdeux}
	\stepcounter{sujet\themeprogrammequn}\stepcounter{sujet\themeprogrammeqdeux}
	\index[v]{Question d'interprétation \typequn!\csname\themeprogrammequn\endcsname!\qun}
	\index[v]{Essai \typeqdeux!\csname\themeprogrammeqdeux\endcsname!\qdeux}
	\index[w]{\csname\themeprogrammequn\endcsname!Question d'interprétation \typequn!\qun}
	\index[w]{\csname\themeprogrammeqdeux\endcsname!Essai \typeqdeux!\qdeux}
	\end{envquestion}
	\clearpage
    }
    %%%%%%%%%%%%%%%%%%%Pour les corrigés%%%%%%%%%%%%%%%%
    \define@cmdkey	[COR]	{corr}	{lieu}			{}
	\define@cmdkey	[COR]	{corr}	{jour}			{}
	\define@cmdkey	[COR]	{corr}	{nosujet}		{}
	\define@cmdkey	[COR]	{corr}	{prenomauteur}	{}
	\define@cmdkey	[COR]	{corr}	{nomauteur}		{}
	\define@cmdkey	[COR]	{corr}	{titre}			{}
	\define@cmdkey	[COR]	{corr}	{type}			{}
	\define@cmdkey	[COR]	{corr}	{liencorr}		{}
	\define@cmdkey	[COR]	{corr}	{annee}			{}
	\define@cmdkey	[COR]	{corr}	{typecorrun}	{}
	\define@cmdkey	[COR]	{corr}	{typecorrdeux}	{}
	\define@cmdkey	[COR]	{corr}	{questcorrun}	{}
	\define@cmdkey	[COR]	{corr}	{questcorrdeux}	{}
	\define@cmdkey	[COR]	{corr}	{corrsautdepage}{}

\newcommand{\corrlieu}{\cmdCOR@corr@lieu}
\newcommand{\corrjour}{\cmdCOR@corr@jour}
\newcommand{\corrnosujet}{\cmdCOR@corr@nosujet}
\newcommand{\corrprenomauteur}{\cmdCOR@corr@prenomauteur}
\newcommand{\corrnomauteur}{\cmdCOR@corr@nomauteur}
\newcommand{\corrtitre}{\cmdCOR@corr@titre}
\newcommand{\corrtype}{\cmdCOR@corr@type}
\newcommand{\corrlien}{\cmdCOR@corr@liencorr}
\newcommand{\corrannee}{\cmdCOR@corr@annee}
\newcommand{\corrtypeun}{\cmdCOR@corr@typecorrun}
\newcommand{\corrtypedeux}{\cmdCOR@corr@typecorrdeux}
\newcommand{\corrquestun}{\cmdCOR@corr@questcorrun}
\newcommand{\corrquestdeux}{\cmdCOR@corr@questcorrdeux}
\newcommand{\corrsautdepage}{\cmdCOR@corr@corrsautdepage}

\presetkeys[COR]{corr}{%
				lieu=Métropole,
				jour={0},
				nosujet=,
				prenomauteur=,
				nomauteur=Anonyme,
				titre={Le Titre},
				type=,
				liencorr=,
				annee=1900,
				typecorrun=,
				typecorrdeux=,
				questcorrun=,
				questcorrdeux=,
				corrsautdepage=,}{}
				
\newcommand{\sujetcorrigen}[3]{
    \setkeys[COR]{corr}{#1}
	\ifthenelse{\equal{\corrnosujet}{}}{%
	\ifthenelse{\equal{\corrtype}{}}{
	\fancyhead[C]{\footnotesize Éléments de correction du jour~\corrjour{}, \corrlieu~(\corrannee)~: \textsc{\corrnomauteur}, \textit{\corrtitre}}
	\subsection{\corrlieu{}~(\corrannee),~Jour~\corrjour}
	}
	{\fancyhead[C]{\footnotesize Éléments de correction du jour~\corrjour{}, \corrlieu~(\corrtype~\corrannee)~: \textsc{\corrnomauteur}, \textit{\corrtitre}}
	\subsection{\corrlieu{}~(\corrtype~\corrannee),~Jour~\corrjour}}
	}
	{%
	\ifthenelse{\equal{\corrtype}{}}
	{\fancyhead[C]{\footnotesize Éléments de correction du sujet~\corrnosujet~du~jour~\corrjour{}, \corrlieu~(\corrannee)~: \textsc{\corrnomauteur}, \textit{\corrtitre}}
	\subsection{\corrlieu~(\corrannee),~Jour~\corrjour,~Sujet~\corrnosujet}}
	{\fancyhead[C]{\footnotesize Éléments de correction du sujet~\corrnosujet~du~jour~\corrjour{}, \corrlieu~(\corrtype~\corrannee)~: \textsc{\corrnomauteur}, \textit{\corrtitre}}	\subsection{\corrlieu~(\corrtype~\corrannee),~Jour~\corrjour,~Sujet~\corrnosujet}}
	}%%%Fin de la mise en page des en-tête et niveaux de titre
\begin{center}
\fbox{
\makebox[\linewidth-3.25cm][c]{
\begin{minipage}{\linewidth-3.25cm}
	\begin{center}
	\noindent {\Large
	\textit{Éléments de correction}}\par
	\ifthenelse{\equal{\corrprenomauteur}{}}
	{%
	\hyperlink{\corrlien}{\textit{Voir le sujet} : \textsc{\corrnomauteur}, \textit{\corrtitre}}\hypertarget{corr\corrlien}{}%
	}
	{%
	{\hyperlink{\corrlien}{\textit{Voir le sujet} : \corrprenomauteur~\textsc{\corrnomauteur}, \textit{\corrtitre}}\hypertarget{corr\corrlien}{}%
	}}
	\end{center}
	\end{minipage}
	}}
\end{center}
\corrint{\corrtypeun}{\corrquestun}\index[c]{Question~d'interprétation~\corrtypeun!\corrquestun{}@{\scriptsize \textbf{\corrlieu,\ \corrannee}~--\  \textsc{\corrnomauteur}, \textit{\corrtitre}}\nopagebreak\\[3pt]\nopagebreak\hspace*{12pt}\corrquestun}

#2

\ifthenelse{\equal{\corrsautdepage}{}}{}{\vfill{~}\clearpage}

\corrref{\corrtypedeux}{\corrquestdeux}\index[c]{Essai~\corrtypedeux!\corrquestdeux{}@{\scriptsize \textbf{\corrlieu,~\corrannee~:}}\nopagebreak\\[3pt]\nopagebreak\hspace*{12pt}\corrquestdeux}

#3
}
 
\makeatother